% Section 8: SYSTEM INTEGRATION AND DATA FLOW

\section{System Integration and Data Flow}

\subsection{Initialization Sequence}
\begin{enumerate}
    \item Serial communication setup (115200 baud)
    \item Logger initialization (rate-limited logging)
    \item LittleFS file system initialization
    \item Settings loading from persistent storage
    \item SensorManager initialization with ADC detection
    \item AutoTuner initialization with PID references
    \item CommandProcessor initialization with all dependencies
    \item PWM channel configuration for valve and analog output
    \item WebHandler initialization with WiFi AP setup
    \item ControlSystem initialization and task creation
    \item TaskManager initialization and FreeRTOS task creation
    \item PID controller configuration and startup
\end{enumerate}

\subsection{Real-time Operation}
The system operates with strict timing requirements:

\begin{table}[H]
\centering
\begin{tabular}{@{}lll@{}}
\toprule
\textbf{Operation} & \textbf{Frequency} & \textbf{Timing Constraint} \\
\midrule
Control Loop & Configurable (default 1000 Hz) & Variable (default 1 ms) \\
Sensor Reading & Configurable (default 1000 Hz) & < 0.5 ms \\
PID Computation & Configurable (default 100 Hz) & 10 ms period \\
Web Data Update & 10 Hz & 100 ms period \\
Serial Commands & On demand & < 10 ms response \\
\bottomrule
\end{tabular}
\caption{System Timing Requirements}
\end{table}

\subsection{Safety Mechanisms}
\begin{itemize}
    \item \textbf{Pressure Limits:} Emergency shutdown on over-pressure via \texttt{isPressureSafe()}
    \item \textbf{Sensor Fallback:} Automatic ADC switching on failure
    \item \textbf{Settings Mutex:} FreeRTOS mutex protects shared settings access (v2.7.3)
    \item \textbf{Watchdog:} FreeRTOS task monitoring
    \item \textbf{PWM Limits:} Valve duty cycle constraints
    \item \textbf{Manual Override:} Emergency manual PWM control
    \item \textbf{Configuration Validation:} Parameter range checking
\end{itemize}

Note: the over-pressure shutdown can currently be defeated by modifying the sensor calibration via the web interface; see Section~\ref{sec:security-compliance}.

\subsection{Settings Mutex and Concurrency}
\label{sec:settings-mutex}

The global \texttt{SettingsHandler settings} object is shared by three execution
contexts: the ControlTask (core 1, high priority, 1000 Hz), the NetworkTask (core 0,
web server), and the Arduino \texttt{loop()} (core 1, serial commands). All three read
and write tunable values such as the setpoint, PID gains, and sensor calibration. To
prevent data races between these tasks, the settings object is protected by a FreeRTOS
mutex. This subsection explains why the mutex is needed and how it is used.

\subsubsection{The Problem: Shared Data and Torn Reads}

On the 32-bit Xtensa LX7 (ESP32-S3), reads and writes of 32-bit values are atomic: they
complete in a single instruction and cannot be interrupted halfway. 64-bit values such
as \texttt{double setpoint} are different. A \texttt{double} occupies 64 bits, so writing
it takes two 32-bit instructions. If the NetworkTask writes the upper 32 bits of a new
setpoint and the ControlTask reads the value before the lower 32 bits are written, the
reader gets half of the old value and half of the new one. This is called a torn read
and produces a value that is neither old nor new, just garbage.

\begin{verbatim}
Core 0 (NetworkTask, writing setpoint = 5.0)   Core 1 (ControlTask, reading setpoint)
---------------------------------------------- --------------------------------------
write upper 32 bits of 5.0 --------------------------------------------------------->
                                                  <-- read upper 32 bits of 5.0 (new)
write lower 32 bits of 5.0                        <-- read lower 32 bits of ??? (old)
\end{verbatim}

For a pressure control system this is dangerous: a torn setpoint or a torn PID gain can
produce a control output that is wildly wrong while still looking valid.

\subsubsection{FreeRTOS Mutex on the ESP32}

The ESP32 Arduino framework uses FreeRTOS. A mutex is a binary semaphore with one extra
property: it records which task owns it, which lets the OS prevent priority inversion
(a low-priority task blocking a high-priority task indefinitely). FreeRTOS provides the
following API:

\begin{lstlisting}[caption=FreeRTOS Mutex API]
SemaphoreHandle_t xSemaphoreCreateMutex();   // Create a mutex
xSemaphoreTake(handle, portMAX_DELAY);       // Lock (block until available)
xSemaphoreGive(handle);                      // Unlock
vSemaphoreDelete(handle);                    // Destroy
\end{lstlisting}

Taking and giving an uncontended mutex takes roughly 1--2\,$\mu$s on the ESP32-S3.

\subsubsection{The Three Access Patterns}

The mutex is a private member of \texttt{SettingsHandler}
(\texttt{SemaphoreHandle\_t \_mutex}), created in the constructor and deleted in the
destructor. Call sites never see the handle. Instead, three methods wrap the lock in
safe, intention-revealing patterns.

\paragraph{Pattern 1: update(fn), writing fields}

\begin{lstlisting}[caption=SettingsHandler::update - Locked Write]
template<typename Fn>
void update(Fn fn) {
    xSemaphoreTake(_mutex, portMAX_DELAY);
    fn(*this);
    xSemaphoreGive(_mutex);
}
\end{lstlisting}

\texttt{update} locks the mutex, calls a lambda with a reference to the settings object,
then unlocks. Callers write one or more fields through the reference:

\begin{lstlisting}[caption=Using update()]
// Single field: filter strength
settings->update([v](SettingsHandler& s){ s.filter_strength = v; });

// Multiple fields atomically: all three PID gains change together
settings->update([&](SettingsHandler& s){
    s.Kp = newKp; s.Ki = newKi; s.Kd = newKd;
});
\end{lstlisting}

The capture clause \texttt{[v]} copies the outer local into the lambda at creation time,
before the mutex is taken. This guarantees the lambda uses the value that was validated
in the outer function, even if another task modifies the original in the tiny gap before
\texttt{update} acquires the lock. \texttt{[\&]} captures by reference and is used when
values must be written back out of the lambda after it runs.

\paragraph{Pattern 2: withLock(fn), protected reads of 64-bit fields}

\begin{lstlisting}[caption=SettingsHandler::withLock - Locked Read]
template<typename Fn>
void withLock(Fn fn) {
    xSemaphoreTake(_mutex, portMAX_DELAY);
    fn();
    xSemaphoreGive(_mutex);
}
\end{lstlisting}

\texttt{withLock} locks the mutex and calls a void lambda. It is used in two places:

\begin{itemize}
    \item \textbf{pid->Compute() in ControlSystem:} the PID library holds a raw pointer
          \texttt{double* mySetpoint} that points directly into \texttt{settings.setpoint}.
          \texttt{Compute()} reads that 64-bit value, so the read must be atomic with
          respect to writes from the NetworkTask:
\begin{lstlisting}[caption=Locked PID Computation]
settings->withLock([&]{ pid->Compute(); });
\end{lstlisting}
    \item \textbf{AutoTuner::process():} the tuner reads the setpoint once per cycle to
          compare it against the measured pressure. The same torn-read risk applies:
\begin{lstlisting}[caption=Locked Setpoint Read]
double sp;
settings->withLock([&]{ sp = settings->setpoint; });
\end{lstlisting}
\end{itemize}

\paragraph{Pattern 3: snapshotTo(out), bulk consistent reads}

\begin{lstlisting}[caption=SettingsHandler::snapshotTo - Atomic Multi-field Read]
void SettingsHandler::snapshotTo(SettingsSnapshot& out) {
    xSemaphoreTake(_mutex, portMAX_DELAY);
    out.Kp = Kp; out.Ki = Ki; out.Kd = Kd;
    out.setpoint = setpoint;
    // ... all other fields ...
    xSemaphoreGive(_mutex);
}
\end{lstlisting}

\texttt{SettingsSnapshot} is a plain struct that mirrors every field of
\texttt{SettingsHandler}. \texttt{snapshotTo} copies all fields in one locked operation,
producing a consistent point-in-time copy. All HTTP GET handlers that render data use it:

\begin{lstlisting}[caption=Using snapshotTo()]
SettingsSnapshot cfg;
settings->snapshotTo(cfg);
// Render the page using cfg.Kp, cfg.setpoint, etc.
// The mutex is released; no lock is held during the slow page send.
\end{lstlisting}

The lock is deliberately not held during the page render: sending an HTTP response can
take milliseconds, and holding the mutex that long would stall the ControlTask every time
a browser loads the page.

\subsubsection{Serialization and Persistence: save() and load()}

\texttt{save()} must read all fields to build a JSON string and then write that string to
flash. Flash I/O is slow (typically 5--20\,ms). If the mutex were held for the whole
operation, the ControlTask would stall for up to 20\,ms on every web-interface change.
The solution is to serialize under the lock and write outside it:

\begin{lstlisting}[caption={SettingsHandler::save - Serialize Under Lock, Write Outside}]
bool SettingsHandler::save() {
    String jsonBuffer;
    xSemaphoreTake(_mutex, portMAX_DELAY);
    JsonDocument doc;
    doc["Kp"] = Kp;
    // ... populate all fields ...
    serializeJson(doc, jsonBuffer);   // fast in-memory operation
    xSemaphoreGive(_mutex);           // release BEFORE the flash write

    File file = LittleFS.open(SETTINGS_FILE_PATH, "w");
    file.print(jsonBuffer);           // slow flash write, no mutex held
    file.close();
    return true;
}
\end{lstlisting}

\texttt{load()} is the mirror image: the file read and JSON parse happen outside the lock,
and only the final field assignments happen under it:

\begin{lstlisting}[caption=SettingsHandler::load - Assign Under Lock]
bool SettingsHandler::load() {
    File file = LittleFS.open(SETTINGS_FILE_PATH, "r");
    JsonDocument doc;
    deserializeJson(doc, file);       // slow I/O and parse, no lock
    file.close();

    xSemaphoreTake(_mutex, portMAX_DELAY);
    Kp = doc["Kp"] | SettingsDefaults::KP;
    Ki = doc["Ki"] | SettingsDefaults::KI;
    // ... all other fields ...
    xSemaphoreGive(_mutex);
    return true;
}
\end{lstlisting}

\subsubsection{Which Fields Need Protection}

Not every field needs the same care. On the ESP32-S3, 32-bit reads and writes are atomic
and can never be torn. Only 64-bit values are at risk.

\begin{table}[H]
\centering
\begin{tabular}{@{}llll@{}}
\toprule
\textbf{Field type} & \textbf{Width} & \textbf{Atomic?} & \textbf{Protection} \\
\midrule
bool, int8\_t & 8 bit & Yes & None \\
float, int, uint32\_t & 32 bit & Yes & None \\
double, int64\_t & 64 bit & No & Must be under lock \\
\bottomrule
\end{tabular}
\caption{Atomicity of Field Types on the ESP32-S3}
\end{table}

In \texttt{SettingsHandler} the only 64-bit fields are \texttt{Kp}, \texttt{Ki},
\texttt{Kd}, and \texttt{setpoint} (all \texttt{double}). All writes go through
\texttt{update()}. The two critical read paths (\texttt{pid->Compute()} and the AutoTuner
setpoint comparison) go through \texttt{withLock()}. Multi-field reads for rendering go
through \texttt{snapshotTo()}. 32-bit fields such as \texttt{filter\_strength} (float),
\texttt{antiWindup} (bool), and \texttt{control\_freq\_hz} (int) are read directly without
locking, because a single-instruction atomic read cannot be torn.

\subsubsection{Summary Diagram}

\begin{verbatim}
NetworkTask (Core 0)              ControlTask (Core 1)
--------------------------        ------------------------------------
Browser sends /set?kp=1.2         pid->Compute() runs at 1000 Hz
  |                                 |
  v                                 v
settings->update([&](s){          settings->withLock([&]{
    s.Kp = 1.2;                       pid->Compute();   // reads setpoint
});                               });
  |                                 |
  Mutex taken for ~2 us             Mutex taken for ~5 us
  (field write only)                (PID arithmetic only)
  |                                 |
settings->save()                  applyHysteresisCompensation()
  |                                 | (reads hysteresis, hystAmount - floats, no lock)
  Mutex taken for ~10 us
  (JSON serialization only)
  |
  Flash write (~10 ms, NO mutex)
\end{verbatim}

The two tasks run at full speed. The only time one waits for the other is during the
brief moments when a field is being written or \texttt{Compute()} is reading the setpoint.

A standalone explanation of the same material is available in
\texttt{sec/2026-08-28\_Mutex\_How\_It\_Works.md}.
